\section{广义相对论}\label{sec:general-relativity}

可以通过保持 Galileo 时空结构的对称 Cartan 联络, 将引力引入 Newton 力学. 因而, 要在狭义相对论中引入引力, 一个自然的想法是寻找保持 Minkowski 内积的对称联络. 为将此想法形式化, 引入下面的定义.

\begin{RiemannDefinition}\label{def:pseudo-riemannian-and-lorentzian-manifolds}
\textbf{伪 Riemann 流形}是一个对 $(M,g)$, 其中 $M$ 是连通的 $n$ 维可微流形, $g$ 是对称非退化的可微 2-张量场 (称 $g$ 为 $M$ 上的\textbf{伪 Riemann 度量}). 伪 Riemann 流形的\textbf{符号差}就是 $g$ 在任意切空间上的符号差. \textbf{Lorentz 流形}是符号差为 $n-2$ 的伪 Riemann 流形 [即对适当选取的标准正交标架, $(g_{\mu\nu})=\operatorname{diag}(-1,1,\ldots,1)$].
\end{RiemannDefinition}

Minkowski 时空 $(\mathbb{R}^4,\langle\cdot,\cdot\rangle)$ 显然是 Lorentz 流形. Levi–Civita 定理的证明用到的是度量的非退化性, 而不是正定性. 因此, 该定理对伪 Riemann 流形仍然成立. 换言之, 给定伪 Riemann 流形 $(M,g)$, 存在唯一的与 $g$ 相容的对称联络 $\nabla$ (由 Koszul 公式给出). 因而, 保持 Minkowski 度量的对称联络只有一个, 即平凡联络 (在 Cartesian 坐标下取所有 Christoffel 符号为零而得到). 注意, 这一联络的测地线是直线, 对应于自由粒子的运动; 特别地, 这些粒子不受任何引力场的作用.

因此, 要通过对称联络引入引力, 就必须考虑更一般的四维 Lorentz 流形, 我们仍将它们称为\textbf{时空}. 不再要求这些时空与 $\mathbb{R}^4$ 微分同胚, 也不要求它们具有惯性坐标卡. 研究这些时空的几何通常称为\textbf{广义相对论}.

每个时空都配有唯一的 Levi–Civita 联络, 从而也有其测地线. 若 $c:I\subset\mathbb{R}\to M$ 为测地线, 则 $\langle\dot c,\dot c\rangle$ 恒定, 因为
\[
  \frac{\mathrm{d}}{\mathrm{d}s}\langle\dot c(s),\dot c(s)\rangle
  =2\left\langle\frac{D\dot c}{\mathrm{d}s}(s),\dot c(s)\right\rangle=0.
\]
按照 $\langle\dot c,\dot c\rangle<0$, $\langle\dot c,\dot c\rangle=0$ 或 $\langle\dot c,\dot c\rangle>0$, 分别称测地线为\textbf{类时}, \textbf{零模}或\textbf{类空} (即按照其切向量是类时, 类空还是零模). 类比 Cartan 联络, 我们用类时测地线表示有质量粒子的自由下落运动. 这保证了等效原理成立. 用零模测地线表示光线的运动.

一般地, 若曲线 $c:I\subset\mathbb{R}\to M$ 满足 $\langle\dot c,\dot c\rangle<0$, 则称其\textbf{类时}. 此时 $c$ 表示非零质量粒子的运动 (除非 $c$ 是测地线, 否则粒子有加速度). 粒子在事件 $c(a)$ 和 $c(b)$ 之间测得的\textbf{固有时}为
\[
  \tau(c)=\int_a^b|\dot c(s)|\,\mathrm{d}s,
\]
其中对任意 $v\in TM$, $|v|=|\langle v,v\rangle|^{1/2}$.

为选出在物理上有意义的时空, 必须施加某种约束. 类比用 Cartan 联络表述 Laplace 方程的方式, 作如下定义.

\begin{RiemannDefinition}\label{def:vacuum-einstein-field-equation}
若 Lorentz 流形 $(M,g)$ 的 Levi–Civita 联络满足 $\operatorname{Ric}=0$, 则称它为 \textbf{Einstein 场方程}的一个\textbf{真空解}.
\end{RiemannDefinition}

一般的 Einstein 场方程为
\[
  \operatorname{Ric}-\frac{S}{2}g=8\pi E,
\]
其中 $S=\sum_{\mu,\nu=0}^3g^{\mu\nu}R_{\mu\nu}$ 为标量曲率, $E$ 为时空中物质的所谓\textbf{能量--动量张量}. 最简单的物质模型是\textbf{无压完美流体}, 由\textbf{静止质量密度函数} $\rho\in C^\infty(M)$ 及单位\textbf{速度向量场} $U\in\mathfrak{X}(M)$ (其积分曲线为流体粒子的运动) 描述. 该物质模型的能量--动量张量为
\[
  E=\rho\,\nu\otimes\nu,
\]
其中 $\nu\in\Omega^1(M)$ 是由度量 $g$ 与 $U$ 相联系的 1-形式.

Einstein 场方程可改写为
\[
  \operatorname{Ric}=8\pi T,
\]
其中
\[
  T:=E-\frac{1}{2}\left(\sum_{\mu,\nu=0}^3g^{\mu\nu}E_{\mu\nu}\right)g
\]
为\textbf{约化能量--动量张量} [参见 Exercise~\ref{ex:general-relativity}(\ref{ex:trace-reversed-einstein-equation})]. 对于无压完美流体, 约化能量--动量张量为
\[
  T=\rho\left(\nu\otimes\nu+\frac{1}{2}g\right),
\]
因而 Einstein 场方程为
\[
  \operatorname{Ric}=4\pi\rho(2\nu\otimes\nu+g)
\]
(将此式与用 Cartan 联络表述的 Poisson 方程比较).

事实表明, 对于适当选取的 $T$, 满足 Einstein 场方程的时空能极其准确地模拟天文现象.

\begin{RiemannExercise}\label{ex:general-relativity}
\begin{enumerate}[label=(\arabic*),ref=\arabic*,leftmargin=*]
\item\label{ex:signature-is-constant-on-pseudo-riemannian-manifold}
证明伪 Riemann 流形 $(M,g)$ 的符号差是良定义的, 即证明 $g_p\in\mathcal{T}^2(T_pM)$ 的符号差不依赖 $p\in M$.
\item\label{ex:trace-reversed-einstein-equation}
证明:
\begin{enumerate}[label=(\alph*),leftmargin=*]
\item Einstein 场方程可以改写为
\[
  \operatorname{Ric}=8\pi T;
\]
\item 对于静止质量密度为 $\rho$, 单位速度 1-形式为 $\nu$ 的无压完美流体, 约化能量--动量张量为
\[
  T=\rho\left(\nu\otimes\nu+\frac{1}{2}g\right).
\]
\end{enumerate}
\item\label{ex:pullback-of-pseudo-riemannian-metric-can-degenerate}
设 $(M,g)$ 为伪 Riemann 流形, $f:N\to M$ 为浸入. 证明 $f^*g$ 不一定是 $N$ 上的伪 Riemann 度量.
\item\label{ex:hyperboloid-model-of-hyperbolic-space}
设 $(M,g)$ 为 $(n+1)$ 维 Minkowski 时空, 即 $M=\mathbb{R}^{n+1}$ 且
\[
  g=-\mathrm{d}x^0\otimes\mathrm{d}x^0+\mathrm{d}x^1\otimes\mathrm{d}x^1
     +\cdots+\mathrm{d}x^n\otimes\mathrm{d}x^n.
\]
设 $i:N\to M$ 为包含映射, 其中
\[
  N:=\{v\in M\mid\langle v,v\rangle=-1\ \text{且 $v^0>0$}\}.
\]
证明 $(N,i^*g)$ 是 $n$ 维双曲空间 $H^n$.
\item\label{ex:fermi-walker-transport}
(\emph{Fermi–Walker 移动}) 设 $c:I\subset\mathbb{R}\to\mathbb{R}^4$ 为 Minkowski 空间中以固有时参数化的类时曲线, $U:=\dot c$ 为单位切向量, $A:=\ddot c$ 为固有加速度. 若
\[
  \frac{DV}{\mathrm{d}\tau}=\langle V,A\rangle U-\langle V,U\rangle A,
\]
则称向量场 $V:I\to\mathbb{R}^4$ 沿 $c$ 作 \textbf{Fermi–Walker 移动}.
\begin{enumerate}[label=(\alph*),leftmargin=*]
\item 证明 $U$ 沿 $c$ 作 Fermi–Walker 移动.
\item 证明若 $V$ 和 $W$ 沿 $c$ 作 Fermi–Walker 移动, 则 $\langle V,W\rangle$ 恒定.
\item 若 $\langle V,U\rangle=0$, 则 $V$ 在 $U$ 处与子流形
\[
  N:=\{v\in\mathbb{R}^4\mid\langle v,v\rangle=-1\ \text{且 $v^0>0$}\}
\]
相切, 而这个子流形与双曲三维空间等距 [参见 Exercise~\ref{ex:general-relativity}(\ref{ex:hyperboloid-model-of-hyperbolic-space})]. 证明在此情形下, $V$ 作 Fermi–Walker 移动当且仅当它沿 $U:I\to N$ 平行移动.
\item 假设 $c$ 描述速率恒为 $v$ 的圆周运动. 设 $V$ 为作 Fermi–Walker 移动的向量场, 与运动平面相切且满足 $\langle V,U\rangle=0$. 求经过一周后 $V$ 转过 (或进动) 的角度.

(\textbf{注}: 可以证明, 自旋粒子的角动量向量沿其运动作 Fermi–Walker 移动, 且与其运动正交; 上述进动已在电子等自旋粒子上观测到, 称为 \textbf{Thomas 进动}.)
\end{enumerate}
\item\label{ex:twin-paradox-on-cylinder}
(\emph{柱面上的双生子佯谬}) 用平移 $\xi(t,x,y,z)=(t,x+8,y,z)$ 生成的离散等距映射群对 Minkowski 时空取商, 得到 Einstein 场方程的一个 (平直) 真空解. 假设地球的运动由直线 $x=y=z=0$ 表示. 再次设 Bob 在两人的 20 岁生日当天以光速的 $80\%$ 沿 $x$ 轴出发, 将自己的双生姐妹 Alice 留在地球上 [参见 Exercise~\ref{ex:special-relativity}(\ref{ex:car-and-garage-paradox})]. 由于空间的拓扑, 两人在 10 年后 (按地球上的测量) 再次相遇, \emph{而 Bob 从未加速}.
\begin{enumerate}[label=(\alph*),leftmargin=*]
\item 求重逢时两人各自的年龄.
\item 从 Bob 的观点看, 是地球远离他. 怎样解释这对双生子的年龄不对称?
\end{enumerate}
\item\label{ex:rotating-frame-in-minkowski-spacetime}
(\emph{旋转参考系})
\begin{enumerate}[label=(\alph*),leftmargin=*]
\item 证明使用 $\mathbb{R}^3$ 中的柱坐标 $(r,\theta,z)$ 后, Minkowski 时空的度量可写为
\[
  g=-\mathrm{d}t\otimes\mathrm{d}t+\mathrm{d}r\otimes\mathrm{d}r
     +r^2\,\mathrm{d}\theta\otimes\mathrm{d}\theta+\mathrm{d}z\otimes\mathrm{d}z.
\]
\item 设 $\omega>0$, 并考虑由 $\theta=\theta'+\omega t$ 给出的坐标变换. 证明在这些坐标下, 度量写为
\begin{align*}
  g&=-(1-\omega^2r^2)\,\mathrm{d}t\otimes\mathrm{d}t
    +\omega r^2\,\mathrm{d}t\otimes\mathrm{d}\theta'
    +\omega r^2\,\mathrm{d}\theta'\otimes\mathrm{d}t\\
   &\quad+\mathrm{d}r\otimes\mathrm{d}r
    +r^2\,\mathrm{d}\theta'\otimes\mathrm{d}\theta'+\mathrm{d}z\otimes\mathrm{d}z.
\end{align*}
\item 证明在区域 $U=\{r<\frac{1}{\omega}\}$ 中, $(r,\theta',z)$ 为常数的坐标曲线是类时曲线, 对应于相对于 $(r,\theta,z)$ 为常数的惯性观察者作刚性旋转的 (加速的) 观察者.
\item 旋转观察者的集合是一个三维光滑流形 $\Sigma$, 具有局部坐标 $(r,\theta',z)$, 且存在自然投影 $\pi:U\to\Sigma$. 如下在 $\Sigma$ 上引入 Riemann 度量 $h$: 若 $v,w\in T_{\pi(p)}\Sigma$, 则
\[
  h(v,w)=g(v^\dagger,w^\dagger),
\]
其中对每个 $u\in T_{\pi(p)}\Sigma$, 向量 $u^\dagger\in T_pU$ 满足
\[
  (\mathrm{d}\pi)_pu^\dagger=u
  \qquad\text{且}\qquad
  g\left(u^\dagger,\left(\frac{\partial}{\partial t}\right)_p\right)=0.
\]
证明 $h$ 是良定义的, 且
\[
  h=\mathrm{d}r\otimes\mathrm{d}r+\frac{r^2}{1-\omega^2r^2}\,
       \mathrm{d}\theta'\otimes\mathrm{d}\theta'+\mathrm{d}z\otimes\mathrm{d}z.
\]
(\textbf{注}: 这是旋转观察者之间作局部距离测量所得的度量; Einstein 利用这一度量有曲率的事实, 论证了在引力的相对论描述中采用非 Euclidean 几何的必要性.)
\item 若曲线 $c:I\subset\mathbb{R}\to U$ 在每一点的切向量 $\dot c$ 都与 $\frac{\partial}{\partial t}$ 正交, 则从旋转观察者的观点看, $c$ 的像由同时发生的事件组成. 证明这等价于要求 $\alpha(\dot c)=0$, 其中
\[
  \alpha=\mathrm{d}t-\frac{\omega r^2}{1-\omega^2r^2}\,\mathrm{d}\theta'.
\]
特别地, 证明一般而言, 沿闭合路径同步旋转观察者的时钟会产生不一致.

(\textbf{注}: 这就是所谓的 \textbf{Sagnac 效应}; 由于地球自转, 同步 GPS 系统地面站中精度很高的原子钟时, 必须考虑这种效应.)
\end{enumerate}
\item\label{ex:static-spacetime-and-fermat-metric}
(\emph{静态时空}) 设 $(\Sigma,h)$ 为三维 Riemann 流形, 考虑由 $M:=\mathbb{R}\times\Sigma$ 及
\[
  g:=-e^{2(\Phi\circ\pi)}\,\mathrm{d}t\otimes\mathrm{d}t+\pi^*h
\]
确定的四维 Lorentz 流形 $(M,g)$, 其中 $t$ 为 $\mathbb{R}$ 上的通常坐标, $\pi:M\to\Sigma$ 为自然投影, $\Phi:\Sigma\to\mathbb{R}$ 为光滑函数.
\begin{enumerate}[label=(\alph*),leftmargin=*]
\item 设 $c:I\subset\mathbb{R}\to M$ 为以固有时参数化的类时测地线, 且 $\gamma:=\pi\circ c$. 证明
\[
  \frac{D\dot\gamma}{\mathrm{d}\tau}=(1+h(\dot\gamma,\dot\gamma))G,
\]
其中 $G=-\operatorname{grad}(\Phi)$ 为由 $h$ 与 $-\mathrm{d}\Phi$ 相联系的向量场, 可看作引力场. 证明这一方程蕴含
\[
  E^2:=(1+h(\dot\gamma,\dot\gamma))e^{2\Phi}
\]
是一个运动常量.
\item 设 $c:I\subset\mathbb{R}\to M$ 为零模测地线, $\widetilde c$ 为用坐标时间 $t$ 对它重新参数化所得的曲线, $\widetilde\gamma:=\pi\circ\widetilde c$. 证明 $\widetilde\gamma$ 是 \textbf{Fermat 度量}
\[
  l:=e^{-2(\Phi\circ\pi)}h
\]
的测地线.

(\textbf{提示}: 利用第~\ref{chap:geometric-mechanics}~章的 Lemma~\ref{lem:levi-civita-under-conformal-change}.)
\item 证明关于 $g$ 的真空 Einstein 场方程等价于
\begin{align*}
  \operatorname{div}G&=h(G,G);\\
  \operatorname{Ric}&=\nabla\mathrm{d}\Phi+\mathrm{d}\Phi\otimes\mathrm{d}\Phi,
\end{align*}
其中 $\operatorname{div}G$ 为 $G$ 的散度, $\operatorname{Ric}$ 和 $\nabla$ 是 $h$ 的 Ricci 曲率和 Levi–Civita 联络, $\nabla\mathrm{d}\Phi$ 为对所有 $X,Y\in\mathfrak{X}(\Sigma)$ 由 $(\nabla\mathrm{d}\Phi)(X,Y):=(\nabla_X\mathrm{d}\Phi)(Y)$ 定义的张量 [参见第~\ref{chap:riemannian-manifolds}~章的 Exercise~\ref{ex:affine-connections}(\ref{ex:covariant-derivative-of-tensors}) 和 Exercise~\ref{ex:levi-civita-connection}(\ref{ex:riemannian-divergence-trace})].
\end{enumerate}
\end{enumerate}
\end{RiemannExercise}
